\section{Setting and Main Contributions}
\label{2405:sec:main:setting}

Given a dataset of $N$ labeled examples $\{\vec z^\nu, y^\nu\}_{\nu \in [N]}$ in $d$ dimensions,\addnote{In this chapter the input covariate is denoted $\vec z$, playing the role of $\vec x$ in the rest of the thesis.} we analyze the learning properties of fully connected two-layer networks with first layer weights $W=\{\vec w_j\}_{j \in [p]}$, second layer $\vec a \in \mathbb{R}^p$, and activation function $\sigma: \mathbb R \to \mathbb R$
\begin{align}
    f(\vec z; W, \vec a) = \frac{1}{p} \sum_{j =1}^p a_j \sigma \left(\langle \vec z, \vec w_j \rangle \right) \label{2405:eq:def_nn}
\end{align}
We consider target functions that are dependent only on few orthogonal relevant directions $k=O_d(1)$ encoded in the target's weight matrix $W^\star = \{\vec w^\star_r\}_{r \in [k]} \in \mathbb{R}^{k\times d}$:
\begin{align}
    y^\nu = f^\star(\vec z^\nu) = h^\star(W^\star \vec z^\nu)
\end{align}
This model for structured data is usually referred to as a \textit{multi-index} model. Our main objective will be to analyze how efficiently two-layer nets adapt to this low-dimensional structure during training. To ensure that all necessary Gaussian expectations exist, we make the following non-restrictive assumption on $\sigma, h^\star$:
\begin{assumption}\label{2405:assump:poly_growth}
    The functions $\sigma$ and $h^\star$ have at most polynomial growth, in the sense that there exists a constant $C > 0$ such that for all $x \in \mathbb{R}, \vec{x} \in \mathbb{R}^k$,
    \begin{equation}\label{2405:eq:poly_growth}
        |\sigma(x)| \leq C(1 + |x|)^C \quad \text{and} \quad |h^\star(\vec{x})| \leq C(1 + \|\vec{x}\|)^C.
    \end{equation} 
\end{assumption}
\paragraph{Training algorithm --}
Our goal is to minimize the following objective:
\begin{equation}
    \cR(W, \vec{a}) = \E_{\vec{z}, y}\left[ \cL(f(\vec{z}, W, \vec{a}), y) \right]
\end{equation}
where $\cL: \R \times \R \to \R$ is a loss function, and the expectation is over the data distribution.

This objective is non-convex, and we do not have direct access to the expectation over the data. A central question in Machine Learning is therefore to quantify the properties of the weights attained at convergence by different algorithms. In this manuscript we consider a specific class of one-pass training routines to minimize the empirical risk. First, we partition the dataset in disjoint mini-batches of size $n_b$, i.e. $\mathcal{D} = \cup_{t \in [T]}\{\vec x^\nu, y^\nu\}_{\nu \in [tn_b, (t+1)n_b]}$ where $T$ is the total number of iterations considered. For each algorithmic step, we update the hidden layer neurons $W=\{\vec w_j\}_{j \in [p]}$ with the program $\mathcal{A}: (\vec w_{j,t}, \gamma, \rho, n_b) \in \mathbb{R}^{d+3} \to \vec w_{t+1,j} \in \mathbb{R}^d$.
\begin{myalgorithm}[Optimizer Main Step]
\label{2405:algo:optimizer_main_step}
 \begin{equation}
 \begin{aligned}
     \vec w_{t+1,j} & =  \mathcal{A}(\vec w_t, \gamma, \rho, n_b) 
     =  \vec w_{t,j} - \frac{\gamma }{n_b}\sum_{\nu \in [n_b]
     }\nabla_{\vec w_{t,j}}
     \cL
    \left(
    f (\vec z^\nu; 
    \Tilde{W}_t(\rho), \vec a_0 ),  y^\nu  \right) \\
   \Tilde{\vec{w}}_{t, j} (\rho) &= \vec{w}_{t, j} - \rho \nabla_{\vec w_{t,j}}
    \cL\left( f(\vec z^\nu; W_t, \vec a_0),  y^\nu \right)
 \end{aligned}
\label{2405:eq:algorithm}     
 \end{equation}
\end{myalgorithm}
This class of algorithms defined by the optimizer step $\mathcal{A}(\cdot , \gamma, \rho, n_b)$, in which we consider the final gradient update after taking a linear combination of the current iterate and its current gradient (noted as $\Tilde{W}_t(\rho)$ in the above), is broadly used in different contexts. Indeed, routines with positive $\rho$ parameters correspond to Extragradient methods (EgD) \citep{korpelevi1976extragradient}, while with negative $\rho$ have been recently used in the context of Sharpness Aware Minimization (SAM) \citep{foret2021sharpness}. For clarity, we present our theoretical results with the above, easily interpretable, Algorithm~\ref{2405:algo:optimizer_main_step}. However, our theoretical claims are valid in a more general setting; we refer to Appendices~\ref{2405:sec:app:proofs}~and~\ref{2405:sec:app:numerics} for additional results on more general optimizer steps. 
% The complete description of the algorithmic routine is given in App.~\ref{2405:sec:app:numerics} (See Algorithm~\ref{2405:algo:full_algorithm}).

As previously stated, the central object of our analysis is the efficiency of the network to adapt to the low-dimensional structure identified by $W^\star$. Therefore, we focus on the learned representations by the first layer weights $W_t$ while keeping the second layer weights $\vec a_t = \vec a_0$ fixed during training. This assumption is favorable to performing the theoretical analysis and is largely used in the theoretical community (e.g., \cite{damian2022neural,ba2022high,bietti2022learning}).

Numerous theoretical efforts have been devoted to understanding the sample complexity needed to learn multi-index targets. However, an important aspect of many algorithmic routines considered is to exploit a clever \textit{warm start} to initialize the iterates \citep{mondelli2018fundamental,luo2019optimal,maillard2020phase,chen2020learning}. We show in this work that such preprocessing of the data is not needed to learn all polynomial multi-index functions in $T=O(d \log d)$ with the routine in Algorithm~\ref{2405:algo:optimizer_main_step}. Instead, we consider a fully agnostic initialization of the network:

\begin{assumption}\label{2405:assump:init}
For each $j \in [p]$, the initial first-layer weight $\vec{w}_{0, j}$ is initialized according to the uniform measure on the sphere:
\[ \vec{w}_{0, j} \overset{i.i.d}{\sim}\operatorname{Unif}(\mathbb{S}^{d-1}). \]
The second-layer weights are drawn i.i.d in $\{-1, 1\}$:
\[ a_{0, j} \overset{i.i.d}{\sim}\operatorname{Unif}(\{-1, 1\}). \]
\end{assumption}

\paragraph{Weak recovery in high dimensions --}
The essential object in our analysis is the evolution of the correlation between the hidden neurons $W = \{\vec w_j\}_{j \in [p]}$ and the target's weights $W^\star$, as a function of the algorithmic time steps. More precisely, we are interested in the number of algorithmic steps needed to achieve an order one correlation with the target's weights $W^\star$, known as \textit{weak recovery}. 
\begin{definition}[Weak recovery]
\label{2405:def:weak_recovery}
    The \textit{target subspace} $V^\star$ is defined as the span of the rows of the target weights $W^\star$:
    \begin{align}
    \label{2405:eq:main:teacher_subspace_def}
        V^\star = \operatorname{span}(\vec w^\star_1, \dots, \vec w^\star_k)
    \end{align}
    We define the weak recovery stopping time for a parameter $\eta \in (0,1)$ independent from $d$ as:
\begin{equation}
    t^{+}_{\eta} = \operatorname{min}\{t \geq 0: \lVert W W^{\star \top} \lVert_F \geq \eta\}
\end{equation}
\end{definition}

\paragraph{From Information Exponents to Generative Exponents --} Recent works have sharply theoretically characterized the weak recovery stopping time for a variety of multi-index targets learned with one-pass SGD. \cite{arous2021online} for the single-index case, then \cite{abbe2022merged} for the multi-index one have unveiled the presence of an Information Exponent $\ell$ characterizing the time complexity needed to weakly recover the target subspace $V^\star$, defined in the single-index case as
\begin{equation}
    \ell = \min \{ k \in \N  :  \Eb{x \sim \cN(0, 1)}{h^\star(x)H_k(x)} \neq 0 \},
\end{equation}
where $H_k$ is the $k$-th Hermite polynomial \citep{odonnell2014analysis}.
Under this framework, the sample complexity required to achieve weak recovery when the link function has Information Exponent $\ell$ is \citep{arous2021online}
\begin{align} \label{2405:eq:gba_time_scaling}
    T(\ell) = 
        \begin{cases}
        &O(d^{\ell -1} ) \qquad \text{if $\ell>2$} \\
        &O(d \log{d} ) \qquad \text{if $\ell=2$} \\
        &O(d) \qquad \text{if $\ell=1$}.
    \end{cases}
\end{align}
These bounds have been improved by \cite{damian2023smoothing} up to order $d^{\ell/2}$; the latter matches the so-called \emph{Correlated Statistical Query} lower bound, which considers queries of the form $\Ea{y \phi(\vec{z})}$.

However, this Information Exponent is not the right measure of complexity for low-dimensional weak recovery. Some iterative algorithms, for instance, are known to perform drastically better\citep{barbier2019optimal,celentano2021high,mondelli2018fundamental,luo2019optimal,maillard2020phase,troiani2024fundamental}. Recently \cite{damian2024computational} introduced a new Generative Exponent $\ell^\star$ (generative exponent) governing the weak recovery time scaling for algorithms in the Statistical Query (SQ) or Low Degree Polynomial (LDP) family. This exponent is defined as
\begin{equation}
    \ell^\star = \min \{ k \in \N  :  \Eb{x \sim \cN(0, 1)}{f(h^\star(x))H_k(x)} \neq 0 \ \text{for some} \ f: \R \to \R \},
\end{equation}
which allows for the application of arbitrary transformations of the labels before performing the queries. Under this definition, the ``hard'' problems with $\ell^\star > 2$ are also impossible to learn for the best first-order algorithm, \emph{Approximate Message Passing} \citep{troiani2024fundamental}. 

We discuss in this work the extension of the generative exponent to the multi-index setting and investigate the dynamics of specific gradient-based programs (Algorithm~\ref{2405:algo:optimizer_main_step}) as a function of the freshly defined hardness exponent $\ell^*$. 
These theoretical findings match the recent one of \cite{kou2024matching} on the computational efficiency of SGD in learning sparse $k-$parity (associated with $\ell^\star = k$).

When the optimal transformation $f$ is known, a simple SGD algorithm on $(\vec{z}, f(y))$ achieves the Generative Exponent lower bound. However, one problem remains: does there exist a class of SGD-like algorithms that can achieve this lower bound agnostically?

\paragraph{Achieving lower bounds with data repetition --}

Although the performance of single-pass SGD is limited by the Information Exponent of $h^\star$, the situation drastically changes when multiple-pass algorithms are considered. Recently, \cite{dandi2024benefits} proved that non-even single-index targets are weakly recovered in $T = O(d)$ when considering extensive batch sizes with multiple pass. This begs the question of how good can a ``reusing'' algorithm be. We answer this question for the class of polynomial link functions:
\addnote{We quote the sharp bound; the additional, not-provably-optimal polylogarithmic factor at generative exponent \(\ge 2\) is suppressed here and carried in full ($d\log(d)^2$) by the formal Theorem~\ref{2405:thm:single_index_recovery}.}
\begin{theorem}[Informal]
    There is a choice of hyperparameters such that if $h^\star$ is a polynomial function, Algorithm~\ref{2405:algo:optimizer_main_step} achieves weak recovery in $O(d\log d)$ samples.
\end{theorem}
This shows in particular that a simple gradient-based algorithm, which only requires seeing the data twice, is optimal for learning any polynomials.

\paragraph{Learning representations --}
Learning relevant features from data is a fundamental property of neural networks. However, much theoretical work has focused on lazy training in two-layer networks \citep{jacot2018neural}, where features are not effectively learned. In this regime, shallow networks behave as kernel machines, which struggle to adapt to low-dimensional structures in data \citep{ghorbani2020when}. Understanding representation learning and corrections to the kernel regime is thus a key challenge in machine learning theory \citep{dudeja2018learning,atanasov2022neural,bietti2022learning,damian2022neural,petrini2022learning}. In this manuscript we provably characterize the number of iterations needed from Algorithm~\ref{2405:algo:optimizer_main_step} to learn the relevant features of the data, i.e., attain weak recovery of the target subspace. Our results provide the scaling of the relevant hyperparameters $(\gamma, \rho)$ with the diverging input dimension $d$. The minibatch size will be fixed to $n_b =1$ in the main body, and we refer to Appendix~\ref{2405:sec:app:numerics} for the extension to larger batch sizes.


\subsection*{Summary of Main Results}
\begin{itemize}
    \item We prove for single-index targets that one-pass gradient-based algorithms surpass the Correlation Statistical Query (CSQ) limits established by the Information Exponent \citep{arous2021online}.  
    \item We unveil that the dynamics of the family of Algorithm~\ref{2405:algo:optimizer_main_step} are governed by the generative exponent \cite{damian2024computational}, with an additional polynomial restriction on the transformation of the output.
    \item We generalize the notion of generative exponent to multi-index targets. As in the single-index case, one-pass gradient-based algorithms can overcome CSQ performance dictated by the Leap Coefficient \citep{abbe2022merged}. Our definition is consistent with the one of \citep{troiani2024fundamental}. 
    \item We prove that all polynomial multi-index functions are learned by Algorithm~\ref{2405:algo:optimizer_main_step} either with total sample complexity $O(d)$ or $O(d \log(d)^2)$ if associated with the presence of a symmetry. 
    \item  The implementation of Algorithm~\ref{2405:algo:optimizer_main_step} does not require pre-processing to initially correlate the estimation with the ground truth, and learns the meaningful representations from data alone. This is in contrast numerous findings (e.g., \cite{mondelli2018fundamental,luo2019optimal,maillard2020phase} for single index targets or \cite{chen2020learning} for multi-index ones).
    \item We characterize the class of hard functions not learned  in (almost) linear sample complexity. We show, however, that such functions can be learned through a hierarchical mechanism that extends the CSQ staircase first developed in \cite{abbe2021staircase} to a different, larger set, of functions.
    \item We validate the formal theoretical claims with detailed numerical illustrations. The code to reproduce representative figures is available in the project \href{https://github.com/IdePHICS/Repetita-Iuvant}{GitHub repository}.
\end{itemize}
